\documentclass[11pt]{article}
\usepackage[utf8]{inputenc}
\usepackage{amsmath,amssymb}
\usepackage{hyperref}
\title{Robust Seismic Tomography Mantle under Distribution Shift}
\author{Anonymous}
\date{}
\begin{document}
\maketitle

\begin{abstract}
This paper studies Seismic Tomography Mantle. Grounded in a real, citable literature \cite{ref1,ref2,ref3,ref4,ref5,ref6,ref7,ref8},
we propose a focused method and report \textbf{illustrative} experiments.
All references below are real and verifiable (OpenAlex/arXiv); the reported
numbers are simulated to illustrate intended behaviour.
\end{abstract}

\section{Introduction}
Seismic Tomography Mantle is an active research area. This work targets a concrete gap identified from
the recent literature and contributes a method together with an evaluation protocol.

\section{Related Work (real literature)}
The following works, retrieved from public bibliographic APIs, frame our study:
\begin{itemize}
\item Müller et al. (2008) — "Age, spreading rates, and spreading asymmetry of the world's ocean crust", Geochemistry Geophysics Geosystems (cited 2082).
\item Wortel et al. (2000) — "Subduction and Slab Detachment in the Mediterranean-Carpathian Region", Science (cited 1618).
\item Hilst et al. (1997) — "Evidence for deep mantle circulation from global tomography", Nature (cited 1275).
\item Huang et al. (2006) — "High‐resolution mantle tomography of China and surrounding regions", Journal of Geophysical Research Atmospheres (cited 1263).
\item Ridolfi et al. (2009) — "Stability and chemical equilibrium of amphibole in calc-alkaline magmas: an overview, new thermobarometric formulations and application to subduction-related volcanoes", Contributions to Mineralogy and Petrology (cited 1204).
\item Zhao et al. (1992) — "Tomographic imaging of <i>P</i> and <i>S</i> wave velocity structure beneath northeastern Japan", Journal of Geophysical Research Atmospheres (cited 1131).
\end{itemize}

\section{Method}
We model distribution shift explicitly and optimise a worst-case objective over an uncertainty set of plausible test distributions. We instantiate this idea for Seismic Tomography Mantle and analyse its cost/quality trade-off.

\section{Experiments (Illustrative)}
\begin{center}
\begin{tabular}{lcc}
System & Primary Metric & Efficiency \\ \hline
Strong baseline & 0.812 & 1.00$\times$ \\
Ablation & 0.828 & 0.92$\times$ \\
\textbf{Ours} & \textbf{0.851} & \textbf{0.71$\times$}
\end{tabular}
\end{center}
\emph{Metrics simulated to illustrate the intended effect; not measured.}

\section{Conclusion}
We connected a real literature to a concrete method for Seismic Tomography Mantle. Future work will
validate the illustrative results with real experiments.

\begin{thebibliography}{9}
\bibitem{ref1} R. Dietmar Müller, M. Sdrolias, Carmen Gaina et al.. Age, spreading rates, and spreading asymmetry of the world's ocean crust. \emph{Geochemistry Geophysics Geosystems}, 2008. DOI:10.1029/2007gc001743
\bibitem{ref2} M. J. R. Wortel, Wim Spakman. Subduction and Slab Detachment in the Mediterranean-Carpathian Region. \emph{Science}, 2000. DOI:10.1126/science.290.5498.1910
\bibitem{ref3} Robert D. van der Hilst, Sri Widiyantoro, E. R. Engdahl. Evidence for deep mantle circulation from global tomography. \emph{Nature}, 1997. DOI:10.1038/386578a0
\bibitem{ref4} Jinli Huang, Dapeng Zhao. High‐resolution mantle tomography of China and surrounding regions. \emph{Journal of Geophysical Research Atmospheres}, 2006. DOI:10.1029/2005jb004066
\bibitem{ref5} Filippo Ridolfi, Alberto Renzulli, Matteo Puerini. Stability and chemical equilibrium of amphibole in calc-alkaline magmas: an overview, new thermobarometric formulations and application to subduction-related volcanoes. \emph{Contributions to Mineralogy and Petrology}, 2009. DOI:10.1007/s00410-009-0465-7
\bibitem{ref6} Dapeng Zhao, Akira Hasegawa, Shigeki Horiuchi. Tomographic imaging of <i>P</i> and <i>S</i> wave velocity structure beneath northeastern Japan. \emph{Journal of Geophysical Research Atmospheres}, 1992. DOI:10.1029/92jb00603
\bibitem{ref7} Philip Wolfe. Convergence Conditions for Ascent Methods. \emph{SIAM Review}, 1969. DOI:10.1137/1011036
\bibitem{ref8} Mark R. Handy, Stefan M. Schmid, Romain Bousquet et al.. Reconciling plate-tectonic reconstructions of Alpine Tethys with the geological–geophysical record of spreading and subduction in the Alps. \emph{Earth-Science Reviews}, 2010. DOI:10.1016/j.earscirev.2010.06.002
\end{thebibliography}
\end{document}
